\title{Byte Latent Transformer: Patches Scale Better Than Tokens}

\begin{abstract}
We present the Byte Latent Transformer (\textsc{BLT}), a novel large language model architecture that operates directly on bytes and, for the first time, reaches parity with token-based LLMs in terms of performance at scale, while also delivering substantial gains in inference speed and resilience. \textsc{BLT} transforms sequences of bytes into patches of variable length, with these patches forming the core computational entities. The patch boundaries are determined according to the predicted entropy of subsequent bytes, ensuring that greater computational focus and model capacity are allocated to segments exhibiting higher complexity. We provide the inaugural flop-matched scaling analysis of byte-level architectures, evaluating models up to 8B parameters and trained on 4T raw bytes. Our empirical findings affirm that it is practical to scale architectures on unprocessed byte streams without relying on a predefined token set. Dynamic selection of larger patches in less complex regions of data enables faster training and inference, with additional qualitative gains observed in reasoning capacity and generalization to rare or complex cases. Taken together, for a given inference budget, \textsc{BLT} exhibits stronger scaling behavior than models relying on tokenization by enabling simultaneous growth in both patch size and model capacity.
\end{abstract}

\section{Introduction}
\label{section:intro}

We present the Byte Latent Transformer~(\textbf{BLT}), an architecture that operates directly on raw byte sequences, entirely eliminating the need for tokenization. BLT achieves parity with established token-based models in terms of performance at scale while offering notable gains in both efficiency and robustness (\S\ref{sec:robustness}). 

Contemporary large language models (llms) are primarily end-to-end systems; however, they still rely on tokenization—an ad hoc pre-processing operation that partitions byte streams into a pre-defined vocabulary of tokens. This process introduces biases into string compression, resulting in challenges such as heightened sensitivity to both the source domain and input modality~\citep{dagangetting}, increased susceptibility to input perturbations~(\S\ref{sec:robustness}), limited understanding of orthographic phenomena~\citep{edman2024cute}, and disparities in multilingual performance~\citep{liang2023xlm,petrov2024language,limisiewicz2024myte}.

Historically, tokenization has played a critical role because training llms directly on raw byte sequences incurs excessive computational expense owing to the considerable sequence lengths involved~\citep{xue2022byt5}. Earlier research addresses this issue by introducing more computationally efficient self-attention mechanisms~\citep{el2020characterbert, clark2022canine} or adopting architectures that eliminate attention altogether~\citep{wang2024mambabyte}~(\S\ref{section:related_work}). Nevertheless, such strategies are mainly beneficial for \textit{small models}. When scaling up, it is the large feed-forward network layers—executed for every byte—that become the primary source of computational demand, rather than the attention component itself.

To optimize compute distribution, we introduce a dynamic, trainable approach for aggregating bytes into \textit{patches}~(\S\ref{section:patching}), as well as a novel model architecture that integrates both byte-level and patch-level information. In contrast to tokenization, {\textsc BLT} does not rely on a static vocabulary for patches. Instead, flexible groupings of bytes are transformed into latent patch representations using efficient, learned encoder and decoder components. Our results demonstrate that this approach enables compute to be allocated \textit{more} effectively than in models that use tokenization.

Tokenization-based LLMs distribute an equal amount of computational resources to every token, resulting in a compromise between computational efficiency and predictive capability. This is because the heuristics used for compression in tokenization are not always aligned with the actual difficulty of predicting certain tokens. The core principle behind our architecture is that computational resources should be adaptively allocated according to the demands of the prediction task. For instance, predicting the completion of a word generally requires less computational effort than generating the initial word of a new sentence, as the latter typically involves higher entropy and more complexity. This principle underpins the structure of {\textsc BLT} (\S\ref{section:architecture}), which employs three transformer modules: two compact byte-level \textit{local models} and a larger, global \textit{latent transformer} (\autoref{fig:arch_high_level}). In order to organize bytes into groups (patches) and to appropriately allocate computational resources, {\textsc BLT} employs segmentation based on the entropy associated with next-byte prediction. This method clusters bytes into context-sensitive patches that maintain a consistent information density.

We conduct the first comprehensive study on scaling byte-level models while keeping the flop count constant, exploring configurations up to 8B parameters and 4T bytes of training data. Our results demonstrate that it is possible to train a large-scale end-to-end model directly from raw bytes, foregoing the need for predetermined vocabulary tokenization. In terms of training performance relative to flop usage,\footnote{The model’s computational requirements are estimated by tallying the total Floating Point OPerations (flops) involved.} {\textsc BLT} achieves parity with Llama 3, yet delivers as much as a 50\% reduction in inference flops~(\S\ref{section:scaling}).

Our analysis further reveals that operating directly on raw byte sequences substantially benefits modeling of rare and diverse data phenomena. {\textsc BLT} outperforms standard tokenizer-driven models in terms of resilience to corrupted input data, and exhibits superior capabilities for character-level processing, as evident in tasks involving orthographic competence, phonological awareness, and low-resource machine translation~(\S\ref{sec:robustness}).

Moreover, utilizing {\textsc BLT} enables simultaneous scaling of both model and patch sizes without surpassing the allocated inference flop budget. Employing larger patch sizes typically lowers the computational load, which permits an expanded global latent transformer, as this component executes fewer times. In experiments designed to hold inference flops constant~(\autoref{fig:fixed_inference_scaling}), we detect notably more favorable scaling patterns than those realized by architectures reliant on tokenization.

To conclude, the key contributions of this work are as follows: 1) We present {\textsc BLT}, an LLM architecture employing byte-level latents that dynamically modulates computational resources, thereby enhancing flop efficiency; 2) We demonstrate that, for models up to 8B parameters, training {\textsc BLT} with an equivalent flop budget achieves parity with Llama 3, while permitting trade-offs—accepting slight declines in evaluation scores for up to a 50\% increase in flop efficiency; 3) Our {\textsc BLT} models introduce a novel approach to scaling large language models, allowing for increased model size without exceeding a fixed inference compute limit; and 4) We show that {\textsc BLT} models exhibit superior robustness to input perturbations and have an improved capacity to capture sub-word input features often overlooked by token-based LLMs. The full training and inference codebase for {\textsc BLT} has been made publicly available at~\url{https://github.com/facebookresearch/blt}.

\section{Patching: From Individual Bytes to Groups of Bytes}
\label{section:patching}

Dividing byte sequences into discrete \textit{patches} enables {\textsc BLT} to flexibly adjust computational resources depending on contextual requirements. As illustrated in Figure~\ref{fig:patching_types}, there are multiple strategies for partitioning bytes into patches. More precisely, a patching function $f_p$ takes a byte sequence $\pmb{x}=\{x_i,|i=1,\ldots n\}$ of length $n$ and segments it into $m$ patches, with $m < n$, represented as $\pmb{p}=\{p_j|j=1,\ldots,m\}$, where each $x_i$ is mapped to either 0 or 1, the latter denoting the commencement of a new patch. The number of main Transformer steps, and thus the primary computational expense in both token-based and patch-based frameworks, hinges on how many patches are produced by a selected patching function in {\textsc BLT}. As such, the average number of bytes per patch—i.e., \textit{patch size}—is the primary determinant of computational cost during both training and evaluation when using a specific patching strategy~(\S\ref{section:flops}).

We proceed by presenting three distinct patching functions: segmenting sequences with a constant byte count per patch~(\S\ref{section:static-patch}), forming patches at whitespace boundaries~(\S\ref{section:white-patch}), and adaptive patching informed by entropy estimations from a lightweight byte language model~(\S\ref{section:dyn-patch}). Additionally, we explore the concept of incremental patching and clarify how this approach differs from conventional tokenization~(\S\ref{section:bpe-patch}).

\subsection{Strided Patching Every K Bytes}
\label{section:static-patch}

One of the simplest methods for aggregating bytes involves organizing them into uniform patches of size $k$, as exemplified by MegaByte~\citep{yu2023megabyte}. This technique, utilizing a constant stride, is convenient to implement for both training and inference stages, offers a direct approach to modifying the average patch size, and thereby facilitates straightforward control over computational requirements. Nevertheless, this approach presents notable limitations. Primarily, the allocation of computation is not adaptive to varying information density: in certain cases, computational resources may be squandered on uninformative regions (e.g., predicting whitespace in code), while areas rich in information, such as mathematical notation, may not receive adequate computation. Additionally, this fixed-size scheme may result in inconsistent segmentation of similar byte sequences—leading, for example, to different splits of identical words based on their positioning.

\subsection{Space Patching}
\label{section:white-patch}

\citet{slagle2024spacebyte} introduces a straightforward but effective enhancement to strided patching: new patches are generated immediately following space-like bytes\footnote{
    Here, space-like bytes are characterized as any byte that does not represent a Latin letter, numerical digit, or UTF-8 continuation byte. Additionally, each patch is required to include at least one byte that is not space-like.}. These boundaries align with common linguistic units in a wide range of languages. In the Space patching approach, a latent transformer operation (requiring increased computational cost) is dedicated to processing each individual word. This strategy leads to consistent patching of words throughout sequences and ensures that additional computational effort is directed toward predictions that tend to be more challenging, which frequently occur after spaces. For instance, when answering the prompt ``Who composed the Magic Flute?\uline{\hspace{.6em}}'', determining the initial byte of the answer poses a considerably greater challenge than predicting subsequent bytes after ``M'' because the first character drastically narrows the set of plausible completions. This renders the rest of the answer ``Mozart'' substantially more predictable. Despite its advantages, space patching is not well-suited for every language or domain and, notably, lacks flexibility in patch size adjustment. In the following, we present a novel patching technique grounded in the observation that the initial bytes of words typically pose the greatest prediction difficulty, while simultaneously enabling straightforward patch size modulation.

\subsection{Entropy Patching: Using Next-Byte Entropies from a Small Byte LM}
\label{section:dyn-patch}

Instead of depending on heuristics like whitespace for segmentation, our method adopts a data-driven strategy to locate points of high uncertainty in next-byte predictions. We propose \textit{entropy patching}, an approach that determines patch boundaries based on estimated entropy values.

We train a small byte-level auto-regressive language model on the training data for {\textsc BLT} and compute next byte entropies under the LM distribution $p_{e}$ over the byte vocabulary $\mathcal{V}$:

\begin{equation}
H(x_i) = \sum_{v \in \mathcal{V}} p_{e}(x_i=v|\pmb{x}_{<i}) \log p_{e}(x_i=v|\pmb{x}_{<i})
\end{equation}

We explore two approaches for detecting patch boundaries using entropies $H(x_i)$. The initial method selects locations where the entropy exceeds a specified global threshold, as demonstrated in~\autoref{fig:patching}. The alternative strategy focuses on points whose entropy is comparatively higher than that of the preceding point. This latter method may also be viewed as detecting points where the generally monotonically decreasing entropy pattern within a patch is disrupted.

\begin{align*}
    \text{Global Constraint}\hspace{1cm}H(x_t)& > \theta_g\\
    \text{Approx. Monotonic Constraint}\hspace{1cm}H(x_t)& - H(x_{t-1}) > \theta_r
\end{align*}

Patch boundaries are determined through an efficient preprocessing stage performed at the point of dataloading. This contrasts with the approach of~\citet{nawrot-etal-2023-efficient}, in which a classifier is trained to identify patch boundaries using entropy-based criteria. Within our experiments~(\S\ref{section:experiments}), we conduct a comparison between these approaches for differentiating low entropy from high entropy bytes.

\subsection{The Byte-Pair Encoding (BPE) Tokenizer and Incremental Patching}
\label{section:incremental-patching}
\label{section:bpe-patch}

Numerous contemporary llms, such as the baseline Llama 3 in this work, utilize subword tokenizers based on bpe~\citep{gage1994new, sennrich-etal-2016-neural}. Throughout this paper, the term ``tokens'' denotes byte sequences selected from a \textit{finite} vocabulary established before training begins, contrasting with ``patches,'' which represent sequences formed dynamically and do not rely on a predetermined vocabulary. An important distinction is that, in the token setting, the model lacks direct interaction with the raw byte-level features.

A significant advancement of {\textsc BLT} over traditional tokenization-based architectures lies in its reconceptualization of the balance between vocabulary size and computational requirements. In conventional llms, expanding the vocabulary leads to, on average, larger tokens, thereby reducing the number of processing steps needed per sequence. However, this also increases the output dimensionality of the model’s final projection layer. As a result, tokenization-based methods face strict limitations when attempting to meaningfully vary either token size or inference computational overhead. For instance, in Llama 3, the mean token size is enlarged from 3.7 to 4.4 bytes, which is achieved by expanding the embedding table size by a factor of four relative to Llama 2.

During generation, {\textsc BLT} must determine at each position in the byte sequence whether a patch boundary has been reached, since this choice triggers whether the Latent Transformer will be utilized for additional computation. Importantly, this decision must be made without referencing future parts of the sequence that have not yet been produced. As a result, patching methods are constrained and cannot rely on knowledge of bytes yet to be generated when segmenting the sequence. More precisely, a patching method $f_p$ possesses the property of incremental patching if it holds that:
$$f_p(\pmb{x}_{<i}) = f_p(\pmb{x})_{<i}$$
bpe fails to qualify as an incremental patching scheme because a given prefix could be tokenized in different ways depending on the subsequent continuation; as such, it does not satisfy the required property above\footnote{Introducing an explicit delimiter to signify patch boundaries can convert bpe into an incremental patching method, though this leads to longer byte sequences.}.

\section{{\textsc BLT} Architecture}
\label{section:architecture}
The {\textsc BLT} system consists of a major global autoregressive language model that processes patch-level representations. Accompanying this are two lightweight local models: one is responsible for transforming byte sequences into patch representations, and the other reconstructs bytes from those representations~(\autoref{fig:arch_high_level}).

\subsection{Latent Global Transformer Model}

\textit{The Latent Global Transformer} refers to an autoregressive transformer architecture $\mathcal{G}$ with $l_{\mathcal{G}}$ layers, which transforms a sequence of latent input patch representations $p_j$ into corresponding output patch representations $o_j$. In this work, subscript $j$ indexes patches, while $i$ indexes bytes. The global model incorporates a block-causal attention mask~\citep{dubey2024llama} that confines each patch's attention to itself and all preceding patches within the document. Notably, this model accounts for the majority of computational workload—both during pre-training and inference—so the timing of its invocation offers a means to dynamically manage and allocate computation across different segments of the input and output, in accordance with their complexity.

\subsection{Local Encoder}
\label{section:local}

\textit{The Local Encoder Model}, represented by $\mathcal{E}$, is designed as a compact transformer model featuring significantly fewer layers ($l_{\mathcal{E}} << l_{\mathcal{G}}$) compared to the global counterpart. Its primary purpose is to transform sequences of input bytes $b_i$ into rich patch representations $p_j$ with high computational efficiency. Distinct from conventional transformer structures, the model integrates a cross-attention mechanism immediately after each transformer block, allowing pooling of byte-level information into patch representations~(\autoref{fig:crossattn}).

Initially, input bytes $b_i$ are mapped to embeddings $x_i$ through a parameter matrix of dimensions $\mathbb{R}^{256 \times h_{\mathcal{E}}}$. These byte embeddings may be further enhanced with hash-based embeddings if needed~(\S\ref{sec:hash-emb}). Subsequently, a sequence of transformer layers and interleaved cross-attention layers~(\S\ref{sec:enc-cross-attn}) iteratively refine these embeddings into patch representations $p_i$, which are later fed to the global transformer module $\mathcal{G}$ for higher-level processing.

Within transformer layers, a \textit{local block causal} attention scheme is utilized: each byte embedding attends to up to $w_{\mathcal{E}}$ preceding byte positions, potentially spanning patch boundaries while ensuring attention is confined to within document limits. The subsections that follow provide elaboration on the embedding strategies and architectural details of the cross-attention blocks.

\subsubsection{Encoder Hash n-gram Embeddings}
\label{sec:ngram-emb}
\label{sec:hash-emb}
An important factor in generating resilient and informative representations at every step $i$ involves leveraging contextual information from previous bytes. Within {\textsc BLT}, this is accomplished by representing not only the single byte $b_i$, but also including it as a constituent of a byte n-gram. At each step $i$, our approach begins with the construction of byte-grams

\begin{align}
    g_{i,n}=\{b_{i-n + 1},\ldots, b_i\}
\end{align}

for each byte position $i$ and for values of $n$ ranging from three to eight.\footnote{
Byte-grams of length $n$ or greater are not used when $i<n$. }

Next, we present hash $n$-gram embeddings, a method that associates each byte $n$-gram to an index within an embedding table $E_{n}^{hash}$—distinct for each $n \in\{3,4,5,6,7,8\}$—by means of a hash function, as previously proposed by~\citet{bai2010learning}. Each such embedding is incorporated into the representation of its corresponding byte. Afterward, the resultant sum undergoes normalization and is used as the input to the local encoder. Specifically, the augmented byte embedding is computed as follows:

\begin{align}
e_i &= x_i + \sum_{n = 3,...,8}E_{n}^{hash}(\text{Hash}(g_{i,n})) \\
\text{where, Hash}(g_{i,n}) &= \text{RollPolyHash}(g_{i,n}) \% |E_{n}^{hash}|
\end{align}

We scale $e_i$ by dividing by the total number of $n$-gram sizes incremented by one, and employ RollPolyHash as specified in Appendix~\ref{appendix:rollpolyhash}. Section~\ref{section:ablations} examines how varying $n$ and embedding table size within $n$-gram hash embeddings impacts scaling law trends under constant flops. Beyond hash $n$-gram embeddings, we have also investigated frequency-based $n$-gram embeddings; further information regarding this approach can be found in Appendix~\ref{appendix:ngrams}.

\subsubsection{Encoder Multi-Headed Cross-Attention}
\label{sec:enc-cross-attn}
Our approach largely adopts the design of the input cross-attention mechanism from the Perceiver architecture~\citep{jaegle2021perceiver}, with a key distinction: instead of employing a fixed pool of latent vectors, the latent states now correspond to adaptable patch representations~(\autoref{fig:crossattn}), and their attention is restricted solely to the bytes comprising their individual patch. Within this framework, each patch $p_j$ is associated with a query vector, which is initialized through a pooling operation over its corresponding byte representations. This pooled result then passes through a linear transformation, $\mathcal{E}_{C} \in \mathbb{R}^{h_{\mathcal{E}} \times (h_{\mathcal{E}} \times U_{\mathcal{E}})}$, where $U_{\mathcal{E}}$ designates the number of encoder cross-attention heads. More concretely, let $f_{\text{bytes}}(p_j)$ represent the sequence of bytes mapped to patch $p_j$; in this case, we compute

\begin{align}
P_{0,j} &= \mathcal{E}_{C}(f_\text{bytes}((p_j)), f~ \text{is a pooling function} \\
P_l &= P_{l-1} + W_o\left(\text{softmax}\left(\frac{QK^T}{\sqrt{d_k}}\right)V\right) \\
\text{where } Q_j &= W_q(P_{l-1,j}), K_i = W_k(h_{l-1,i}), V_i = W_v(h_{l-1,i}) \\
h_l &= \text{Encoder-Transformer-Layer}_l(h_{l-1})
\end{align}

Here, $P \in \mathbb{R}^{n_p \times h_{\mathcal{G}}}$ denotes the set of $n_p$ patch representations provided as input to the global model, with initialization performed by aggregating the byte-level embeddings $e_i$ that constitute each patch $p_j$. The matrices $W_q$, $W_k$, $W_v$, and $W_o$ serve as the learned projections for queries, keys, values, and output, respectively; keys and values are obtained by projecting the byte-wise hidden states $h_i$ from the previous layer (or $e_i$ at the first layer). Our masking scheme is adapted to the patching context, restricting each query $Q_j$ to only access those keys and values derived from bytes within the same patch $j$. Owing to the deployment of multi-head attention across $Q$, $K$, and $V$, and since patch representations generally have a higher dimensionality ($h_{\mathcal{G}}$) than $h_{\mathcal{E}}$, we preserve $P_l$ as several heads of size $h_{\mathcal{E}}$ during cross-attention, concatenating them subsequently to recover $h_{\mathcal{G}}$ dimensions. Furthermore, pre-LayerNorm is applied to queries, keys, and values, and this cross-attention component omits any positional encodings. A residual connection is also incorporated around the cross-attention layer.

\subsection{Local Decoder} 

Analogous to the architecture of the local encoder, the local decoder $\mathcal{D}$ is also implemented as a compact transformer model, characterized by a depth of $l_{\mathcal{D}}$ layers, where $l_{\mathcal{D}} << l_{\mathcal{G}}$. Its purpose is to transform the global patch representations $o_j$ back into the original raw byte sequence, denoted as $y_i$. During generation, the local decoder outputs each byte in sequence, conditioned on all previously reconstructed bytes. To accomplish this, it utilizes the hidden states outputted by the local encoder corresponding to the sequence of bytes. The structure of the decoder comprises $l_{\mathcal{D}}$ iterations of cross-attention modules followed by transformer blocks. In this setup, each cross-attention module is positioned before its associated transformer layer, enabling the model to construct byte-level representations from the input patch embeddings prior to the subsequent byte-wise processing performed by the transformer layer.

\subsubsection{Decoder Multi-headed Cross-Attention} 

In the decoder cross-attention, the roles of the queries and key/values are interchanged i.e. the byte-representations are now the queries, and the patch representations are now the key/values. The initial byte-representations for the cross-attention are initialized as the byte embeddings from the last encoder layer i.e. $h_{l_{\mathcal{E}}}$. The subsequent byte-representations for layer $l$, $d_{l,i}$ are computed as:

\begin{align}
D_0 &= h_{l_{\mathcal{E}}} \\
B_l &= D_{l-1} + W_o\left(\text{softmax}\left(\frac{QK^T}{\sqrt{d_k}}\right)V\right), \\
  &\text{where } Q_i = W_q(d_{l-1,i}), K_i = W_k(\mathcal{D}_C(o_j)), V_i = W_v(\mathcal{D}_C(o_j)) \\
  D_l &= \text{Decoder-Transformer-layer}_l(B_l)
\end{align}

Here, $W_k$ and $W_v$ denote the projection matrices for keys and values, which are applied to the final patch representations $o_j$ produced by the global model via the linear transformation and split operation $\mathcal{D}_C$. The query projection matrix, $W_q$, acts on the byte representations $d_{l-1}$ from the prior decoder transformer layer, or on $h_{l_{\mathcal{E}}}$ when considering the initial layer. The output projection is accomplished using $W_o$. Consequently, $B \in \mathbb{R}^{h_{\mathcal{D}} \times n_b}$, with $n_b$ representing the total number of output bytes. Subsequent decoder representations $D_l$ are generated through a decoder transformer layer that operates on the results of the cross-attention module, $B$. In a manner analogous to the cross-attention in the local encoder, we adopt multi-head attention, employ pre LayerNorms, omit positional embeddings, and incorporate a residual connection encircling the cross-attention sublayer.

\section{Experimental Setup}
\label{section:experiments}
We systematically construct controlled experimental protocols to compare {\textsc BLT} and tokenization-based models, explicitly ensuring that {\textsc BLT} does not benefit from extended sequence contexts, in order to avoid any unfair advantage in our evaluations.

\subsection{Pre-training Datasets} For all model configurations examined in this study, pre-training is conducted on two main datasets: 1) The Llama 2 dataset~\citep{touvron2023llama}, containing 2 trillion tokens derived from a diverse array of publicly accessible sources, subjected to comprehensive cleaning and filtering processes to enhance data integrity; and 2) {\textsc BLT}-1T, a newly constructed dataset encompassing 1 trillion tokens sourced from multiple public repositories, inclusive of a selection of pre-training samples from Datacomp-LM's released data~\citep{li2024datacomp}. The Llama 2 dataset facilitates scaling law analyses regarding the ideal token count, as previously established by~\cite{dubey2024llama}, with the goal of optimizing {\textsc BLT}'s architecture. The {\textsc BLT}-1T dataset supports complete pre-training procedures to benchmark against Llama 3 performance on downstream evaluation tasks. Importantly, both datasets are free from data originating from Meta products or services. Additionally, for control experiments involving tokenizer-based models, we employ the Llama 3 tokenizer, featuring a 128K token vocabulary, as it demonstrated superior baseline performance compared to the Llama 2 tokenizer in our tests.

\subsection{Entropy Model} For our experiments, the entropy model is implemented as a byte-level language model, leveraging the same training data as the comprehensive {\textsc BLT} model. Unless specified otherwise, we employ a transformer architecture comprising 100M parameters, spread across 14 layers, with a hidden dimension of 512 and a sliding window attention that covers 512 bytes. All other hyperparameters align with those utilized in our local and global transformer configurations. Additional experiments involving variations in model size, receptive field, and alternative architectures are detailed in \autoref{section:ablations}. Notably, when the model’s receptive field is sufficiently narrow, the learned entropy model can be compactly represented as a lookup table.

\subsection{Entropy Threshold and Equalizing Context Length} For models that utilize entropy-based patching, we determine a patching threshold to reach a targeted average \textit{patch size} over the pretraining data mix. In the case of {\textsc BLT}, as opposed to conventional tokenization, the \textit{patch size} is a tunable parameter, which notably influences the effective context size processed by the model. To standardize the average context length and prevent models with larger patch sizes from obtaining an undue benefit, we ensure the total number of bytes per batch is constant on average. This necessitates decreasing the sequence length for models with increased patch sizes. For Llama 2 data, we employ an 8k byte context, whereas for the {\textsc BLT}-1T dataset, we extend the average context length to 16k bytes while preserving the batch size at an average of 16M bytes.

Although the mean batch size remains unchanged, dynamic patching strategies result in varying byte-to-patch proportions when assembling data batches. To optimize efficiency, our approach to {\textsc BLT} training compiles batches using a fixed number of patches, which eliminates the need for additional padding within the computationally intensive latent transformer. Consequently, all batches are standardized to contain an identical count of patches. To prevent memory usage surges caused by exceptionally large patch sequences during training, we apply padding and, when necessary, truncation to the byte sequences, setting limits at 12k and 24k bytes for the Llama 2 and {\textsc BLT}-1T datasets, respectively.

\subsection{Entropy Model Context}
\label{section:entropy-context}
Our empirical observations indicate that applying entropy patching tends to produce increasingly larger patches, particularly in structured datasets such as multiple choice tasks, where content tends to be highly repetitive (see \autoref{fig:mmlu_patching} for a patching illustration using an MMLU sample). These fluctuations can be attributed to diminished entropy arising from recurrent elements within the entropy model context. To address this, during the extensive run of {\textsc BLT}-Entropy with a patch size of 4.5, we refresh the entropy context by introducing new line characters and employ the approximate monotonicity constraint, as it is less susceptible to "entropy drift" associated with variable context lengths. This adjustment solely alters the method for entropy calculation; our process for determining the entropy threshold remains unchanged.

\subsection{FLOPs Estimation} 
\label{section:flops}

Our approach closely mirrors the methodology outlined in Chinchilla~\citep{hoffmann2022training} for calculating transformer flops, which encompasses operations within feed-forward layers, qkvo projections in the self-attention mechanism, as well as both attention computation and output projection. A key distinction in our analysis is the assumption that the input embedding layer employs an efficient lookup rather than dense matrix multiplication, effectively resulting in zero flop consumption for this component. Consistent with established practice, we approximate the computational cost of the backward pass to be double that of the forward pass.

To compute flops \textit{per byte} for {\textsc BLT} models, we add up the flops for the local encoder transformer, the global latent transformer, and the local decoder transformer, together with the cross attention blocks in the encoder and the decoder:

\begin{align}
    \text{FL}_{\text{{\textsc BLT}}} &= \text{Transf. FL}(h_{\mathcal{G}}, l_{\mathcal{G}}, m=n_{ctx}/n_p, V=0) \div n_p\\
   &\qquad + \text{Transf. FL}(h_{\mathcal{E}}, l_{\mathcal{E}}, m=w_{\mathcal{E}}, V=0) \\
   &\qquad + \text{Transf. FL}(h_{\mathcal{D}}, l_{\mathcal{D}}, m=w_{\mathcal{D}}, V=256) \\ 
    &\qquad + \text{Cross Attn. FL}(h_{\mathcal{E}}, l_{\mathcal{E}}, m=n_p, r=n_p/k) \cdot k/n_p \\ 
   &\qquad + \text{Cross Attn. FL}(h_{\mathcal{D}}, l_{\mathcal{D}}, m=k, r=k/n_p)
\end{align}

Here, $n_{ctx}$ denotes the length of the input sequence measured in bytes, and $n_p$ specifies the patch size. The parameter $r$ represents the proportion between queries and key/values, while $k$ indicates the factor relating patch-dimension to byte-dimension, essentially signifying the number of local model partitions merged to form the comprehensive model representation (as illustrated with $k=2$ in \autoref{fig:crossattn}). The vocabulary size $V$ pertains to the output projection and is utilized exclusively in the local decoder. The length of the context $m$ in the attention mechanism depends on whether the operation targets the byte-level or the patch-level sequence. As a result, we adapt the calculation of attention flops for each particular component. A detailed breakdown of the equations used to compute Transformer-FLOPs and Cross-Attention FLOPs can be found in Appendix~\ref{section:flops-appendix}.

\subsection{Bits-Per-Byte Estimation} Perplexity only makes sense in the context of a fixed tokenizer as it is a measure of the uncertainty for each token. When comparing byte and token-level models, following previous work~\citep{xue2022byt5, yu2023megabyte, wang2024mambabyte}, we instead report Bits-Per-Byte (BPB), a tokenizer independent version of perplexity. Specifically:

\begin{align}
    \text{BPB}(x)&= \frac{\mathcal{L}_{CE}(\pmb{x})}{\ln(2)\cdot n_{\text{bytes}}}
\end{align}

Here, the cumulative cross-entropy loss quantifying the uncertainty associated with the data $\pmb{x}$ is divided by the product of the number of bytes comprising $\pmb{x}$ and a fixed constant, serving to normalize the result.

\subsection{Transformer Architecture Hyperparameters} The configuration of all transformer components within {\textsc BLT}, encompassing both local and global modules, generally aligns with the design of Llama~3~\citep{dubey2024llama}. Specifically, feed-forward sublayers utilize the SwiGLU activation~\citep{swiglu}, while rotary positional embeddings (RoPE)~\citep{RoPe}, parameterized by $\theta=500000$~\citep{xiong2024effective}, are applied exclusively to self-attention layers. RMSNorm~\citep{RMSNorm} serves as the normalization strategy across all layers. For self-attention layers employing static attention patterns such as \textit{block causal} or \textit{fixed-window block causal}, we implement Flash Attention~\citep{dao2022flashattention} with a window length of 512 tokens for fixed-width masking. In cross-attention modules, due to the presence of adaptive, patch-dependent attention masks, Flex Attention\footnote{\url{https://pytorch.org/blog/flexattention}} is leveraged for fused execution, resulting in considerable improvements in training efficiency.

\subsection{{\textsc BLT}-Specific Hyperparameters} In order to evaluate the performance of the {\textsc BLT} models, we perform experiments in two primary domains: analyzing scaling behaviors and conducting assessments on downstream tasks. For these analyses, we select models of varying scales: 400M, 1B, 2B, 4B, and 8B parameters. Details regarding the architectural hyperparameters for each scale can be found in Appendix~Table~\ref{tab:arch}. 

For the initialization of queries within the first cross-attention layer of the local encoder, we apply a max-pooling strategy. All {\textsc BLT} model variants utilize $500,000$ hashes generated via a single hash function, employing n-grams with lengths spanning from 3 to 8. The learning rate is uniformly set to $4e-4$ across all scales. This particular learning rate is selected based on hyperparameter sweeps conducted between $1e-3$ and $1e-4$ at the 400M and 1B scales, indicating equivalence in optimality for both token and {\textsc BLT} models. For scaling investigations on Llama-2 datasets, batch sizes align with recommendations in~\cite{dubey2024llama}, converted to an equivalent number of bytes when necessary. We utilize the AdamW optimizer~\citep{adamw} with coefficients $\beta_1$ and $\beta_2$ fixed at 0.9 and 0.95, respectively, and $\epsilon=10^{-8}$. The optimization process involves a linear warm-up phase spanning 2000 steps, followed by cosine decay of the learning rate to zero. Furthermore, we incorporate a weight decay term of 0.1, and apply global gradient clipping capped at 1.0.

\section{Scaling Trends}
\label{section:scaling}

We offer a comprehensive overview of how byte-level models scale, providing insights valuable for guiding the continued expansion of {\textsc BLT} models. Our examination of scaling seeks to overcome the shortcomings found in earlier analyses of byte-level approaches by: (a) contrasting trends within a compute-efficient training paradigm, (b) constructing parallel 8B-sized models using substantial training datasets—reaching as high as 1T tokens or 4T bytes—and evaluating their efficacy on a variety of downstream benchmarks, and (c) analyzing how scaling behaves under constant inference cost conditions. In a subsequent section, we delve into particular benefits that emerge from directly processing byte sequences.

\subsection{Parameter Matched Compute Optimal Scaling Trends}
\label{section:parameter-matched-scaling}
We leverage the Llama 2 dataset to train several \textit{compute-optimal} bpe and {\textsc BLT} models at four model capacities, specifically spanning from 1B to 8B parameters. To evaluate, we chart the total training flops in relation to language modeling accuracy, measured over a representative fraction of the training mixture. The bpe models employ the parameter-to-data scaling ratio recommended by Llama~3~\citep{dubey2024llama}, ensuring that, under a fixed computation budget, model performance is maximized~\citep{hoffmann2022training}. This forms a strong and theoretically motivated baseline for our experiments. Correspondingly, for each bpe model, a {\textsc BLT} model is also trained on identical data; these use Latent Transformers configured to mirror the size and structure of their bpe Transformer counterparts.

As shown in~\autoref{fig:scaling-compute-opt} (right), {\textsc BLT} models consistently perform on par with or better than their bpe counterparts, a pattern that persists across increases in model size and computational effort. To our knowledge, {\textsc BLT} represents the first byte-level Transformer framework to reach scaling behavior comparable to that of BPE-based architectures under compute-optimal conditions. These findings support our hypothesis that the optimal relationship between the number of parameters and training compute determined for bpe extends to {\textsc BLT}, or at the very least deviates minimally.

Both enhancements in architecture and the use of dynamic patching are essential to achieve scaling behaviors similar to those observed with BPE. In ~\autoref{fig:scaling-compute-opt} (left), we present a comparison between models that employ space patching and Llama 3. We represent SpaceByte \citep{slagle2024spacebyte} as the {\textsc BLT} model with space-patching, but exclude n-gram embeddings and cross-attention modules for this estimation. While SpaceByte demonstrates performance gains over Megabyte, there remains a significant gap relative to Llama 3. \autoref{fig:scaling-compute-opt} (right) showcases the incremental benefits accrued from architectural modifications in conjunction with dynamic patching. At scale, {\textsc BLT} models achieve results comparable to state-of-the-art models that utilize tokenization, such as Llama 3.

Additionally, our results indicate that for models relying on tokenizers, utilizing the Llama-3 tokenizer leads to better performance compared to employing the Llama-2 tokenizer, even when both are trained on identical datasets.

In summary, when operating with much larger patch sizes, the {\textsc BLT} architecture exhibits scaling behavior that falls between that of Llama 2 and Llama 3. While the bpe tokenizers employed by Llama 2 and Llama 3 have mean token lengths of 3.7 and 4.4 bytes respectively, {\textsc BLT} demonstrates comparable scaling characteristics even as the average patch size increases to 6 or 8 bytes. Because inference flops are inversely related to the average patch size, adopting an 8-byte patch size would yield almost a 50\% reduction in inference compute. Furthermore, models configured with larger patch sizes tend to show improved performance as both model capacity and dataset scale increase. For example, {\textsc BLT} initialized with a patch size of 8 lags behind bpe-tokenized Llama 2 at the 1B parameter mark but surpasses the bpe baseline by the 7B parameter scale. This pattern indicates the potential for superior performance with large patch sizes at even greater model scales and suggests the viability of experimenting with even larger patches as model size and training computation budgets expand.

\subsection{Beyond Compute Optimal Task Evaluations}
\label{section:evals}
To further investigate scaling behavior, we extend training of an 8B {\textsc BLT} model past the compute-optimal threshold using the {\textsc BLT}-1T dataset, which is both more extensive and of higher quality. We then evaluate model performance across a range of established benchmarks in classification and generation. For this assessment, we include tasks spanning common sense reasoning, world knowledge, and code generation:

\paragraph{Classification tasks} consist of ARC-Easy (0-shot)~\citep{Eval_ARC}, Arc-Challenge (0-shot)~\citep{Eval_ARC}, HellaSwag (0-shot)~\citep{Eval_hellaswag}, PIQA (0-shot)~\citep{Eval_piqa}, and MMLU (5-shot)~\citep{hendrycks2020measuring}. Accuracy is computed by evaluating the likelihood of different response options using prompt-based scoring, and the mean accuracy across tasks is reported.

\paragraph{Coding related generation tasks:}  The model’s coding capabilities are gauged using pass@1 metrics on MBPP (3-shot)~\citep{Eval_MBPP} and HumanEval (0-shot)~\citep{Eval_HumanEval}, examining proficiency in generating Python code.

In \autoref{tab:evals}, we present a comparative analysis of three models, each trained on the {\textsc BLT}-1T dataset: a model utilizing the Llama 3 tokenizer with a bpe vocabulary,\footnote{We select the Llama 3 tokenizer, featuring a 128k vocabulary, due to its superior performance over the Llama 2 tokenizer, which has a 32k vocabulary.} and two distinct variants of the {\textsc BLT} model. The first variant adopts a space-patching strategy ({\textsc BLT}-Space), while the second leverages an entropy-based patching mechanism ({\textsc BLT}-Entropy). The latter model is further guided by an approximate monotonicity constraint, and we periodically reset its context with new line tokens, as detailed in~\autoref{section:entropy-context}. All three approaches are trained under an equalized computational budget (in flops). For the {\textsc BLT}-Entropy model, we also introduce an inference-time modification by lowering the entropy threshold from 0.6 to 0.1. This adjustment leads to improved task performance, albeit requiring additional inference steps.

The {\textsc BLT}-Entropy model achieves superior results compared to the Llama 3 model on 4 out of 7 evaluated tasks, even though both models are trained with an equivalent data size in bytes. This enhanced performance can likely be attributed to two factors: (1) a more efficient utilization of training computation enabled by dynamic patching, and (2) the explicit modeling of information at the byte level rather than relying on tokenization.

Conversely, {\textsc BLT}-Space lags behind the Llama 3 tokenizer in every task except one; however, it offers notable savings in inference flops due to its increased mean patch size of 6 bytes. In contrast, models utilizing bpe and entropy-patching approaches possess an average patch size near 4.5 bytes when evaluated on the mixed training data. Given equivalent training resources, the model with the larger patch size processes 30\% more data than the alternatives, which may result in {\textsc BLT} deviating further from the compute-optimal regime.

\subsection{Patches Scale Better Than Tokens} The {\textsc BLT} architecture enables concurrent scaling of both model and patch size without altering the training or inference flop allocation and while using a fixed quantity of training data. Patch-based models, unlike their token-based, fixed-vocabulary counterparts, allow for patch sizes to be increased without bound, providing a pathway to bypass traditional efficiency limitations as outlined in Section~\ref{section:bpe-patch}. Utilizing larger patches reduces overall computational demands, allowing the conserved resources to be redirected toward enlarging the global latent transformer, which benefits from being invoked at a reduced frequency.

We perform a fixed inference scaling analysis to examine whether using larger models that process fewer steps over bigger patches yields superior performance compared to smaller models with more steps. Initiating from 400m and 3.6B parameter models utilizing the Llama 2 tokenizer, we identify flop-matched counterparts employing the Llama 3 tokenizer, as well as {\textsc BLT}-Entropy variants with mean patch sizes of 6 and 8 bytes within the training datamix (model specifications are outlined in~\autoref{tab:fixed-inf-params}). For models with a patch size of 8, we implement 3 encoder layers, in contrast to just 1 for smaller patch configurations. Each model is trained under a range of training flop constraints.

\autoref{fig:fixed_inference_scaling} demonstrates that {\textsc BLT} models exhibit superior scaling properties compared to tokenization-based systems across both categories of inference flops. For smaller training budgets, BPE models initially outperform; however, as training investment increases past the compute-optimal point, {\textsc BLT} models soon take the lead. In practice, allocating additional resources during the pretraining phase—despite being a single, up-front cost—can yield models that significantly outperform others for a constant inference budget. This pattern is well illustrated by the 8B model class exemplified by Llama 3.1, which was pretrained on data quantities two orders of magnitude larger than what would be optimal if only compute constraints were considered for that model size.

The threshold at which {\textsc BLT} begins to surpass token-based models moves marginally closer to the compute-optimal regime as the models scale up in FLOPs (decreasing from 3 times to 2.5 times the compute-optimal allocation). Moreover, within the higher-FLOP regime, models employing a patch size of 8 display a sharper improvement curve, allowing them to outpace their counterparts at an earlier stage. As outlined in Section~\ref{section:parameter-matched-scaling}, expanding patch sizes tends to enhance performance relative to BPE-based models at greater model capacities. This phenomenon can be partially explained by the reduced proportion of overall FLOPs allocated to the byte-level Encoder and Decoder, which do not scale their computational demands as rapidly as the Latent Transformer does. For example, when scaling up total model parameters by a factor of 20, from 400M to 8B, the local parameters assigned specifically to {\textsc BLT} increase only about twofold. This distinction matters because larger patch sizes modify only the computational load from the patch Latent Transformer, leaving the byte-level module requirements unchanged. Consequently, this explains why the {\textsc BLT}-Entropy with patch size 8 shifted from occupying 1.6x to 1.7x the size of the Llama 2 model as the scale of the models increased.

To summarize, our investigation into patch-length scaling reveals that the {\textsc BLT} patch-based architecture exhibits superior scaling properties when both the patch size and model size are enlarged concurrently. These favorable scaling behaviors not only persist but may further strengthen as model size increases.

\section{Byte Modeling Improves Robustness} We further evaluate the robustness of {\textsc BLT}, comparing it to token-based models that do not utilize explicit byte-level representations, and introduce a method for augmenting pretrained token-based models with byte-level processing.
\label{sec:robustness}

\subsection{Character-Level Tasks}

One of the initial incentives for developing byte-level models was their inherent resilience to noise at the byte level within inputs, as well as their capability to recognize and process the internal structure of tokens—an area where traditional tokenizer-based approaches often encounter challenges. To assess these properties, we conduct supplementary evaluations using benchmarks designed to test both the robustness of models to input noise and their sensitivity to character-level information, spanning English, multiple languages, numerical digits, and phonetic representations. The findings from these evaluations are summarized in Table~\ref{tab:char_tasks}.

\paragraph{Noisy Data} To evaluate robustness, we generate perturbed versions of the benchmark classification datasets outlined in Section~\ref{section:evals}, allowing for direct comparison between {\textsc BLT} and traditional tokenizer-based architectures. Five unique character-level perturbation methods are used to modify the input text: (a) \textit{AntSpeak}: The complete text is rendered in uppercase, with characters separated by spaces. (b) \textit{Drop}: Removes 10\% of text characters at random. (c) \textit{RandomCase}: Each character is randomly converted to either uppercase or lowercase with equal probability across the text. (d) \textit{Repeat}: 20\% of characters are duplicated, with each repetition limited to at most four occurrences. (e) \textit{UpperCase}: The entirety of the text is shifted to uppercase letters. During the evaluation phase, these noise procedures are applied independently to prompts, completions, or both, with performance measured by averaging over these scenarios. The corresponding results for noised HellaSwag~\citep{Eval_hellaswag} are displayed in Table~\ref{tab:char_tasks}. Our findings show that {\textsc BLT} consistently demonstrates greater robustness than tokenizer-dependent models, outperforming the baseline model (trained on identical data) by an average margin of 8 points, and even exceeding the performance of Llama 3.1, despite that model's exposure to a significantly larger training corpus.

\paragraph{Phonology - Grapheme-to-Phoneme (G2P)} We evaluate the effectiveness of {\textsc BLT} in translating sequences of graphemes—alphabetic characters forming a word—into their corresponding phonemic representations, which capture the word's pronunciation. The outcomes of the G2P evaluation, shown in Table \ref{tab:char_tasks}, utilize a 5-shot setting from the Phonology Bench~\citep{suvarna2024phonologybench}. Our findings indicate that {\textsc BLT} surpasses the performance of the baseline Llama 3 1T tokenizer-based model on this task.

\paragraph{CUTE}
To evaluate fine-grained character-level comprehension, we test {\textsc BLT} using the CUTE benchmark~\citep{edman2024cute}, which consists of multiple tasks grouped into three principal domains: compositional understanding, orthographic similarity, and sequence manipulation capabilities. The benchmark poses a notable difficulty for traditional tokenizer-based models; these models generally retain token spelling information but have limited facility in leveraging it for dynamic text manipulation. As presented in Table~\ref{tab:char_tasks}, {\textsc BLT}-Entropy surpasses both variants of BPE Llama 3 models by over 25 points on the CUTE benchmark. Notably, our architecture achieves near-perfect accuracy (99.9\%) on both spelling-oriented manipulation tasks, exhibiting remarkable adeptness at handling character-level modifications. This significant performance gain is evident despite {\textsc BLT} being trained with merely one-sixteenth the data used for Llama 3.1, suggesting BPE-based approaches have inherent difficulty acquiring robust character-level competencies. Selected cases depicted in Figure~\ref{fig:cute_examples} highlight instances where Llama 3 with BPE tokenization falls short, whereas {\textsc BLT} handles them successfully. The two scenarios—word deletion and word insertion—are exceptions in which the BPE-based model slightly outperforms ours; nevertheless, this class of word-level editing is arguably more complex for byte-level frameworks, yet the margin of difference remains modest. Furthermore, transitioning from character-based manipulation to word-based edits may be a more tractable challenge than the reverse. Consistent evaluation protocols were maintained across all tasks using unaltered Huggingface prompts. It is worth noting that BPE models may see improved results through targeted prompt engineering.

\paragraph{Low Resource Machine Translation} We assess the performance of {\textsc BLT} for both directions of translation involving twenty-one low-resource languages, spanning six major language families and employing diverse scripts, as represented in the FLORES-101 benchmark~\citep{goyal2022flores}. SentencePiece BLEU scores are presented in Table~\ref{tab:flores}. Our findings indicate that {\textsc BLT} surpasses a counterpart trained using the Llama 3 tokenizer, showing an overall improvement of 2 BLEU points when translating into English and a 0.5-point gain when translating out of English. While translation quality in widely used language pairs is similar or marginally improved over Llama 3, the advantages of {\textsc BLT} are especially pronounced across many lower-resource language pairs. These results highlight the strengths of byte-based modeling in capturing long-tail byte sequence regularities.

\subsection{Training {\textsc BLT} from Llama 3}
\label{sec:distillation}
We investigate a methodology where {\textsc BLT} models benefit from pre-existing, pre-trained tokenizer-based models to enhance both the efficiency and speed of training. This is accomplished by initializing the global transformer parameters in {\textsc BLT} using those from a pre-trained Llama 3.1 model. During training, we fine-tune the global transformer weights with a learning rate set to one-tenth of that applied to the local encoder and decoder, maintaining the optimal number of training steps used for Llama 3. Table~\ref{tab:distill} compares this approach with a baseline {\textsc BLT}. Results clearly show that {\textsc BLT} initialized from Llama 3.1 exceeds the performance of both baseline {\textsc BLT} and Llama 3 models, even when identical amounts of computational resources (flops) are expended. Furthermore, when contrasted with our {\textsc BLT}-Entropy variant in Table~\ref{tab:evals}, which was trained using a much more extensive dataset of 1T tokens, {\textsc BLT} initialized from Llama 3.1 continues to demonstrate improved performance on the MMLU benchmark. This suggests that the proposed method provides an effective strategy for considerably decreasing the computational cost required for training.

This approach can be interpreted as enabling tokenizer-based architectures to function without tokenizers, thereby transforming a pre-trained LLaMA 3.1 model into a {\textsc BLT} architecture. For a thorough evaluation, we report results for the original LLaMA 3.1 model (trained on 15T tokens) in Table \ref{tab:distill} and contrast its performance with that of the {\textsc BLT} obtained from LLaMA 3. Our results indicate that although there is only a minor decrease in performance on MMLU and HumanEval, more pronounced drops are observed on several other benchmarks. These findings highlight the need for additional research to more effectively utilize the pre-trained model's capabilities, particularly through improved optimization of data mixtures and other hyperparameters.

\section{Ablations and Discussion}
\label{section:ablations}
Here, we present a series of ablation studies aimed at validating the design decisions underlying {\textsc BLT}, as well as evaluating the patching strategy and chosen hyper-parameters for the {\textsc BLT} 8B parameter configuration, which was trained using the {\textsc BLT}-1T dataset.

\paragraph{Entropy Model Hyper-parameters} To investigate how entropy model capacity and context window size affect scaling, we conduct experiments using byte-level entropy transformer architectures, varying the number of parameters from 1 million up to 100 million, and adjusting context windows across the range of 64 to 512 tokens. The corresponding bits-per-byte (bpb) against training compute (FLOPs) scaling behavior is illustrated with the $400m$ and $1b$ {\textsc BLT} models, which are trained on the Llama-2 dataset (see \autoref{fig:entropy_ablation}). The results demonstrate that both increased model size and longer context contribute to improved scaling, although additional gains become marginal for models larger than 50 million parameters.

\paragraph{Types of Patching}

We conduct ablation studies on the four patching methods outlined in Section~\ref{section:patching}: (1) Strided Patching employing strides of 4 and 6, (2) Whitespace-based Patching, (3) Patching using BPE Tokenizer aligned with the Llama 3 tokenizer, and (4) Entropy-driven patching leveraging a lightweight byte-level language model.

Although dynamic patching shortens the apparent sequence length, we adjust the sequence length parameter to ensure that each patching method processes comparable context lengths. Consequently, every model is exposed to an equivalent number of bytes per sequence, both during training and inference, thereby eliminating potential confounds related to context size differences. The findings of these ablation studies are illustrated in Figure~\ref{fig:scaling-compute-opt}. Notably, all alternative patching strategies achieve superior performance compared to static patching, with space patching demonstrating results that are nearly equivalent to those of dynamic entropy-based patching.

\autoref{tab:evals_ablation} reports on the performance of {\textsc BLT} models, contrasting traditional tokenizer-based approaches, space patching, and entropy-based patching, with all models trained on the Llama 2 dataset for the recommended number of steps~\citep{dubey2024llama}. While space patching offers a straightforward alternative that bypasses the need for on-the-fly entropy model computation during training, our results indicate that the advantages identified for entropy-based patching with respect to scaling patterns~(Section~\ref{section:scaling}) are also evident in a range of downstream benchmark assessments.\footnote{The results for space patching derive from previous experiments lacking cross-attention, though comparable trends are seen when cross-attention is employed.}

\paragraph{Cross-Attention} 
In~\autoref{tab:crossattn_ablations_1b}, we investigate the impact of introducing cross-attention modules at different stages within the encoder and decoder components of {\textsc BLT}. Regarding the encoder cross-attention, we evaluate three strategies for initializing the queries: 1) assigning the same learned embedding across all global states, 2) utilizing a hash-based embedding of the patch bytes, and 3) employing a pooled version of the encoder’s hidden states corresponding to the patch bytes at that particular encoder layer.

Our results indicate that incorporating cross-attention within the \textit{decoder} yields the greatest benefit. When applied in the encoder, cross-attention offers modest gains, contingent on initializing queries through pooling. Moreover, we observe that cross-attention proves especially advantageous for the Common-Crawl dataset, with the positive impact being more pronounced as patch sizes increase.

\paragraph{n-gram Hash Embeddings} 

We investigate the effects of varying n-gram hash embedding vocabulary sizes—specifically 0, 100K, 200K, and 400K—as detailed in Table~\ref{tab:ngram_ablations_1b}. Our results demonstrate consistent improvements from hash embeddings across all evaluated domains, with Wikipedia and Github showing particularly notable gains (an observed difference of 0.04 bpb versus 0.01 bpb after 15k steps using the 8B model). At the 8B parameter scale, a reduction from 500K to 300K hash entries only marginally impacted performance, with changes around ~0.001 bpb at the 15k step mark. These findings suggest that hash embeddings are essential for enhancing {\textsc BLT} models to a level comparable with tokenizer-based approaches; nonetheless, performance increases plateau beyond 300K hashes. Moreover, our analysis indicates that the performance benefits offered by hash embeddings are largely additive with those obtained from cross-attention, as each method yields improvements on distinct datasets.

\paragraph{Local Model Hyperparamaters} Table \ref{tab:local_layers} presents an ablation study over different configurations for the number of layers in both the local encoder and decoder. Utilizing hash n-gram embeddings, {\textsc BLT} demonstrates strong performance even with a minimal encoder architecture—specifically, a single-layer encoder—combined with a more substantial decoder.

\section{Related Work}
\label{section:related_work}

\textbf{Character-Level RNNs:} The task of Character Language Modeling has remained central to the development of neural approaches~\citep{sutskever2011generating,mikolov2012subword,graves2013generating}, primarily because such models inherently handle novel or out-of-vocabulary terms without relying on back-off techniques. \cite{kim2016character} introduced a model that solely utilizes characters for input processing, leveraging convolutional and highway networks integrated into LSTM-based RNNs. This architecture not only achieved results on par with contemporary RNN-based language models for English but also demonstrated superior outcomes for morphologically rich languages, highlighting a notable benefit of character-level LLMs. Similarly, \cite{kenter2018byte} implemented byte-level LSTM models in machine comprehension tasks, where they consistently outperformed word-level models in languages such as Turkish and Russian that have complex morphology. In a related effort, \cite{zhang2015character} applied character-based convolutional architectures to classification challenges, showing notable advantages over word-level approaches in particular tasks. Further enhancing this approach, \citet{chung2019hierarchical} proposed hierarchical LSTM architectures equipped with boundary detectors at each level to uncover implicit hierarchical structures within text, thereby advancing the effectiveness of character-level language modeling. Alternatively, ByteNet introduced by \cite{kalchbrenner2016neural} adopted convolutional layers at the character level, moving away from attention mechanisms for tasks like machine translation.

\textbf{Character-Level Transformers:} Transformer architectures leveraging attention mechanisms~\citep{vaswani2017attention} and subword tokenization techniques~\citep{sennrich-etal-2016-neural} have notably advanced the performance of neural language models across various tasks and benchmarks. Nonetheless, selecting word or subword tokens inherently introduces a level of abstraction that constrains model inductive bias. Seeking to unify the benefits of transformer models with earlier advancements in character-based language modeling, \citet{al2019character} explored extremely deep transformer architectures and incorporated auxiliary objectives during training. Their transformer-based character models achieved superior results compared to traditional LSTM-based character language models. Despite these improvements, a substantial performance discrepancy remained when compared with word-level large language models. Similarly, GPT-2~\citep{radfordgpt2} reported that, on expansive datasets such as the 1 Billion Word Benchmark, language models operating at the byte level lagged behind their word-level counterparts.

\cite{Choe2019BridgingTG} showed that transformer-based language models working at the byte level can achieve better performance compared to subword-level LLMs with similar parameter counts; however, these models demand significantly greater computational resources and longer training durations. Likewise, \cite{el2020characterbert} introduced CharFormer, a variant of BERT that constructs word representations by applying convolutional layers over character embeddings, and reported notable gains within the medical domain, although these results come at the expense of higher computation costs. The CANINE model developed by \cite{clark2022canine} comprises 150M parameters and processes character sequences directly through an encoder-only architecture. Its core design involves a deep transformer stack, analogous to our global model, and integrates a local transformer with strided convolutions to reduce the length of input character sequences, surpassing the performance of the similarly sized token-based encoder-only model (mBERT) on various multilingual downstream tasks. ByT5~\citep{xue2022byt5} investigated byte-level encoder-decoder architectures, intentionally foregoing patching mechanisms. Although these models showed greater resilience to input noise and rivaled tokenizer-based models with only a quarter of the data, omitting patching resulted in costly attention operations across all bytes, making computation highly intensive. Modeling input sequences at the byte level, rather than with subword units, leads to increased sequence lengths, introducing new challenges to the scalability of such models. More recently, leveraging the Mamba Architecture~\citep{gu2023mamba}—which supports maintaining a fixed-size memory state over long input contexts—\cite{wang2024mambabyte} trained a byte-level Mamba model (again without patching), and were able to achieve better bits-per-byte performance compared to byte-level transformer counterparts in flop-controlled experiments at the 350M parameter scale, across multiple benchmark datasets.

\textbf{Patching-based approaches:} Employing patching strategies can significantly decrease the substantial flop requirements typically associated with byte-level LLMs, all while possibly preserving their effectiveness. Several studies have reported promising results in limited settings regarding both model size and volume of training data. \cite{nawrot-etal-2022-hierarchical} conduct experiments utilizing static patching that incorporates both downsampling and upsampling, culminating in the development of the hourglass transformer. This model achieves superior performance compared to existing byte-level baselines at the 150M parameter scale. Building on this, \cite{nawrot-etal-2023-efficient} propose enhancements via dynamic patching techniques, such as a boundary-predictor trained end-to-end, a boundary-predictor guided by specific tokenizers, and an entropy-based patching approach resembling {\textsc BLT}. Their findings show that these innovations can exceed the capabilities of standard transformers on language modeling tasks, specifically for models with 40M parameters trained on roughly 400M tokens. In related work, \cite{lester2024training} explore the use of arithmetic coding for sequence compression, obtaining compression levels surpassing those of BPE. Through an equal-info windows strategy, their approach yields better results than byte-level baselines on language modeling benchmarks, though it remains less effective than subword-based baselines.

Our study is heavily influenced by MegaByte~\citep{yu2023megabyte}, a causal LLM structured as a decoder-only model. MegaByte utilizes a fixed static approach for patching, where input bytes are grouped into patches and their representations concatenated; on the decoder side, a local model reconstructs the original byte sequence from these patches. Their results indicate that, at the 1B parameter scale and across datasets containing approximately 400B bytes, MegaByte achieves performance on par with models reliant on tokenization. In the course of our own experiments, we systematically compare against MegaByte and observe that the static patching strategy falls short when compared with compute-efficient, tokenizer-based models operating under identical FLOPs constraints. Our results illustrate that {\textsc BLT} effectively narrows this performance deficit. Further, \cite{slagle2024spacebyte} arrive at similar conclusions regarding the limitations of MegaByte, proposing modifications such as patching on whitespace and analogous byte patterns, as well as the incorporation of a local encoder model. Their findings show improvements relative to transformer models based on tokenization, particularly within specific domains such as Github and arXiv at the 1B parameter level. We include empirical evaluations of this extended model and highlight that additional architectural enhancements are essential for byte-level models to not only scale more effectively but also to reach parity with leading token-based architectures like Llama 3.

\section{Limitations and Future Work}

For the architectural decisions in this study, we align our model training durations with the optimal number of steps identified for Llama~3~\citep{dubey2024llama}. It should be noted, however, that these scaling laws were originally derived from experiments involving BPE-level transformers, and their direct application to {\textsc BLT} might not yield the best ratios between dataset and parameter size. Determining specific scaling laws for {\textsc BLT}, which may demonstrate even more advantageous scaling properties for our architecture, remains an open avenue for future investigation. Furthermore, the majority of our experiments have been limited to models containing up to 1B parameters, so the ideal choices regarding architecture could potentially differ as we extend to 8B parameters or more—a regime where further gains in performance may be realized.

Current transformer repositories and code frameworks are optimized primarily for tokenizer-based transformer models. Although we conduct theoretically flop-matched experiments and utilize some efficient components (like FlexAttention) to support layers that differ from the standard transformer configuration, our implementations might still lag behind tokenizer-based models concerning wall-clock performance and could potentially be improved with additional optimizations.

Although {\textsc BLT} employs a separately pre-trained entropy model for patching, an intriguing avenue for future investigation involves integrating patching model training into a fully end-to-end pipeline. As described in Section \ref{sec:distillation}, we conduct preliminary experiments suggesting that it is feasible to ``byte-ify'' tokenizer-based architectures such as Llama 3—trained on datasets exceeding 10T tokens—by transferring their weights to the global transformer and freezing them. Advancing research along this path may yield techniques that preserve the advantages of bytefying, and even surpass the performance of the original tokenizer-based systems, all without necessitating training from the beginning.

\section{Conclusion}
\label{section:conclusion}

In this work, we introduce the Byte Latent Transformer (\textbf{BLT}), an innovative model architecture that moves beyond the standard reliance on static vocabulary-based tokenization employed in conventional large language models. \textsc{BLT} utilizes a flexible, trainable approach for aggregating bytes into variable-sized patches, thereby dynamically distributing computational effort in accordance with the underlying data's complexity. This shift leads to notable gains in both computational efficiency and model robustness. Through a comprehensive scaling analysis, we show that {\textsc BLT} attains parity with tokenization-dependent models such as Llama 3, even at model sizes up to 8B parameters and with 4T bytes of data, while also enabling up to a 50\% decrease in inference floating-point operations with only minimal drops in evaluation performance. Additionally, {\textsc BLT} introduces an innovative pathway for model scaling by allowing simultaneous adjustment of both model and patch sizes, all while maintaining a constant inference budget. Such adaptability proves particularly beneficial in computational environments reflective of real-world constraints. By directly processing input at the byte level, {\textsc BLT} not only augments the handling of data exhibiting long-tail distributions but also strengthens resilience against noisy data and improves sub-word structure modeling. Taken together, our findings establish {\textsc BLT} as a compelling alternative to established tokenization-centric methods, delivering a versatile, efficient, and robust foundation for next-generation language models.

\end{document}